\section{Step 12 Use the reflected transition for table construction}\label{crc-s12}

Apply the same bit transition to the entry remainder in \texttt{t0}. This is the operation being precomputed, so its polynomial and shift direction must agree with the byte kernel. A separately specified table oracle will check the initializer rather than trusting agreement between two student routines.

\Needspace{12\baselineskip}
\subsection{Write and inspect}
Replace the marked temporary instructions in \texttt{i\_step}. Retain surrounding labels and markers.

\MARSCodeLabel{MIPS assembly}
\begin{minted}[samepage=true]{mars-idea}
    andi $t2, $t0, 1
    dsrl $t0, $t0, 1
    beq $t2, $zero, init_no_feedback
    nop
    xor $t0, $t0, $t3
\end{minted}

Use \texttt{s12-1}. Prepare the exact memory inputs in a separate source file and pause at \texttt{kernel\_done}, before the caller stores the result.

\subsection{Expected observation}
Entry zero remains zero after this first transition.

\Needspace{11\baselineskip}
\subsection{Check your understanding}
Why check table words independently?

\begin{itemize}
\item \textbf{A} A matching wrong initializer and updater can conceal a shared error.
\item \textbf{B} Only the first table word matters.
\item \textbf{C} A table makes the polynomial irrelevant.
\end{itemize}
